% TicketManager — User stories Sprint 1 (aligné sur l'implémentation réelle)
% Exporter les .drawio en PNG dans le dossier img/ (voir noms ci-dessous).

% ============================================================
% US01 — S'authentifier
% ============================================================
\subsubsection{Description textuelle du cas d'utilisation « S'authentifier »}

\begin{table}[H]
\centering
\footnotesize
\renewcommand{\arraystretch}{0.8}
\setlength{\tabcolsep}{4pt}
\begin{tabular}{|p{3cm}|p{9cm}|}
\hline
\textbf{Élément} & \textbf{Description} \\
\hline
Cas d'utilisation & S'authentifier \\
\hline
Acteurs & Client, Technicien, Administrateur \\
\hline
Précondition & L'acteur possède un compte valide (email + mot de passe). Le backend NestJS et MongoDB sont disponibles. \\
\hline
Postcondition & Un JWT (\texttt{access\_token}) est émis, stocké côté client (\texttt{localStorage}) et l'acteur est redirigé vers son espace selon son rôle (\texttt{/dashboard}, \texttt{/technician} ou \texttt{/admin}). \\
\hline
Scénario nominal &
1. L'acteur accède à \texttt{/login}.\\
& 2. Il saisit email et mot de passe, puis clique sur « Se connecter ».\\
& 3. Le frontend appelle \texttt{POST /auth/login}.\\
& 4. \texttt{AuthService} vérifie les identifiants (bcrypt) et \texttt{JwtService} génère un JWT (24 h).\\
& 5. Le frontend enregistre le token et le rôle, puis redirige l'utilisateur.\\
\hline
Scénarios alternatifs &
• Identifiants incorrects (401 — « Email ou mot de passe incorrect »).\\
& • Email invalide ou mot de passe $<$ 6 caractères (400).\\
& • Backend injoignable (message d'erreur réseau).\\
& • Accès à un espace protégé sans token (redirection \texttt{/login}).\\
\hline
\end{tabular}
\caption{Description textuelle du cas d'utilisation : S'authentifier}
\label{tab:cu-auth}
\end{table}

La figure~\ref{fig:seq-auth} illustre ce scénario.

\begin{figure}[htbp]
    \centering
    \includegraphics[width=0.85\textwidth]{img/seq-authentification.png}
    \caption{Diagramme de séquence -- Authentification JWT}
    \label{fig:seq-auth}
\end{figure}

\subsubsection{Développement et intégration de l'interface d'authentification}
L'interface de connexion (\texttt{/login}) permet la saisie sécurisée des identifiants avec affichage des erreurs. Après validation, le système délivre un JWT, stocke le token et le rôle côté client, puis redirige automatiquement l'utilisateur vers son espace selon son profil RBAC.


% ============================================================
% US02 — Création d'un ticket avec pièces jointes
% ============================================================
\subsubsection{Description textuelle du cas d'utilisation « Créer un ticket avec pièces jointes »}

\begin{table}[H]
\centering
\footnotesize
\renewcommand{\arraystretch}{0.8}
\setlength{\tabcolsep}{4pt}
\begin{tabular}{|p{3cm}|p{9cm}|}
\hline
\textbf{Élément} & \textbf{Description} \\
\hline
Cas d'utilisation & Créer un ticket avec pièces jointes \\
\hline
Acteurs & Client (secondaire : Administrateur notifié) \\
\hline
Précondition & Le client est authentifié (JWT, rôle \textit{client}). Cloudinary est configuré pour l'upload. \\
\hline
Postcondition & Un ticket est créé en base avec statut \texttt{open}, historique \texttt{created}, URLs Cloudinary dans \texttt{attachments[]}. Les administrateurs sont notifiés (WebSocket + base). \\
\hline
Scénario nominal &
1. Le client accède à \texttt{/dashboard/tickets/new}.\\
& 2. Il renseigne titre, description, catégorie et priorité.\\
& 3. Il joint éventuellement des fichiers (glisser-déposer ou sélection).\\
& 4. À la soumission : \texttt{POST /tickets/attachments/upload} puis \texttt{POST /tickets}.\\
& 5. Le ticket est persisté dans MongoDB ; notifications admin émises.\\
& 6. Redirection vers \texttt{/dashboard/tickets}.\\
\hline
Scénarios alternatifs &
• Titre vide (400 — validation).\\
& • Fichier $>$ 10 Mo ou type non autorisé (JPEG, PNG, GIF, WebP, PDF).\\
& • Échec upload Cloudinary ou backend injoignable.\\
& • Session expirée (401 — redirection \texttt{/login}).\\
\hline
\end{tabular}
\caption{Description textuelle du cas d'utilisation : Créer un ticket avec pièces jointes}
\label{tab:cu-create-ticket}
\end{table}

La figure~\ref{fig:seq-create-ticket} illustre ce scénario.

\begin{figure}[htbp]
    \centering
    \includegraphics[width=0.85\textwidth]{img/seq-creation-ticket-pieces-jointes.png}
    \caption{Diagramme de séquence -- Création d'un ticket avec pièces jointes}
    \label{fig:seq-create-ticket}
\end{figure}

\subsubsection{Développement et intégration de l'interface de création de ticket}
L'interface propose un formulaire (titre, description, catégorie, priorité) et une zone de dépôt de fichiers. Les pièces jointes sont uploadées sur Cloudinary avant l'enregistrement du ticket ; le client est redirigé vers la liste de ses tickets après confirmation.


% ============================================================
% US03 — Consultation, filtrage et modification des tickets personnels
% ============================================================
\subsubsection{Description textuelle du cas d'utilisation « Consulter, filtrer et modifier ses tickets »}

\begin{table}[H]
\centering
\footnotesize
\renewcommand{\arraystretch}{0.8}
\setlength{\tabcolsep}{4pt}
\begin{tabular}{|p{3cm}|p{9cm}|}
\hline
\textbf{Élément} & \textbf{Description} \\
\hline
Cas d'utilisation & Consulter, filtrer et modifier ses tickets personnels \\
\hline
Acteurs & Client \\
\hline
Précondition & Le client est authentifié (JWT, rôle \textit{client}). \\
\hline
Postcondition & Le client visualise sa liste paginée filtrée ; en cas de modification réussie, le ticket \texttt{open} est mis à jour avec entrée \texttt{history[updated]}. \\
\hline
Scénario nominal &
1. Le client accède à \texttt{/dashboard/tickets}.\\
& 2. Il applique des filtres : statut, priorité, catégorie, dates, recherche texte.\\
& 3. \texttt{GET /tickets} retourne uniquement \textbf{ses} tickets (\texttt{findAllByUser}).\\
& 4. Il consulte le détail via \texttt{TicketDetailModal}.\\
& 5. Pour un ticket \texttt{open}, il modifie via \texttt{/dashboard/tickets/[id]/edit} puis \texttt{PATCH /tickets/:id}.\\
\hline
Scénarios alternatifs &
• Aucun ticket pour les filtres sélectionnés (liste vide).\\
& • Modification refusée si statut $\neq$ \texttt{open} (403).\\
& • Erreur chargement ou session expirée (401).\\
\hline
\end{tabular}
\caption{Description textuelle du cas d'utilisation : Consulter et gérer ses tickets}
\label{tab:cu-consult-ticket}
\end{table}

La figure~\ref{fig:seq-consult-ticket} illustre ce scénario.

\begin{figure}[htbp]
    \centering
    \includegraphics[width=0.85\textwidth]{img/seq-consultation-tickets-client.png}
    \caption{Diagramme de séquence -- Consultation, filtrage et modification des tickets}
    \label{fig:seq-consult-ticket}
\end{figure}

\subsubsection{Développement et intégration de l'interface de gestion des tickets client}
L'interface « Mes Tickets » offre une liste paginée avec filtres combinés et rafraîchissement temps réel (WebSocket). Le client peut modifier titre, description, catégorie, priorité et pièces jointes tant que le ticket reste au statut « Ouvert ».


% ============================================================
% US04 — Assignation d'un ticket à un technicien
% ============================================================
\subsubsection{Description textuelle du cas d'utilisation « Assigner un ticket à un technicien »}

\begin{table}[H]
\centering
\footnotesize
\renewcommand{\arraystretch}{0.8}
\setlength{\tabcolsep}{4pt}
\begin{tabular}{|p{3cm}|p{9cm}|}
\hline
\textbf{Élément} & \textbf{Description} \\
\hline
Cas d'utilisation & Assigner un ticket à un technicien \\
\hline
Acteurs & Administrateur (secondaires : Technicien, Client notifiés) \\
\hline
Précondition & L'administrateur est authentifié (rôle \textit{admin}). Le ticket existe. Au moins un technicien est enregistré. \\
\hline
Postcondition & Le champ \texttt{assignedTo} pointe vers le technicien choisi ; historique \texttt{assigned\_changed}. Notifications WebSocket + persistantes au technicien et au client. \textbf{Le statut n'est pas modifié automatiquement.} \\
\hline
Scénario nominal &
1. L'administrateur ouvre la liste des tickets (\texttt{/admin}).\\
& 2. Il ouvre \texttt{AssignTicketModal} et charge \texttt{GET /users?role=technician}.\\
& 3. Il sélectionne un technicien et confirme.\\
& 4. \texttt{PATCH /tickets/:id} avec \texttt{\{ assignedTo: technicianId \}}.\\
& 5. \texttt{handleAdminUpdate()} met à jour le ticket et notifie les acteurs concernés.\\
\hline
Scénarios alternatifs &
• Réassignation ou retrait d'affectation (\texttt{assignedTo: null}).\\
& • Utilisateur non admin (403).\\
& • Erreur serveur ou session expirée (401).\\
\hline
\end{tabular}
\caption{Description textuelle du cas d'utilisation : Assigner un ticket}
\label{tab:cu-assign-ticket}
\end{table}

La figure~\ref{fig:seq-assign-ticket} illustre ce scénario.

\begin{figure}[htbp]
    \centering
    \includegraphics[width=0.85\textwidth]{img/seq-assignation-ticket.png}
    \caption{Diagramme de séquence -- Assignation d'un ticket à un technicien}
    \label{fig:seq-assign-ticket}
\end{figure}

\subsubsection{Développement et intégration de l'interface d'assignation de ticket}
La modal \texttt{AssignTicketModal} permet de rechercher et sélectionner un technicien. L'assignation met à jour le ticket via l'API et déclenche des notifications au technicien (« Ticket affecté ») et au client (« Technicien affecté »).


% ============================================================
% US05 — Traitement d'un ticket et mise à jour du statut
% ============================================================
\subsubsection{Description textuelle du cas d'utilisation « Traiter un ticket et mettre à jour le statut »}

\begin{table}[H]
\centering
\footnotesize
\renewcommand{\arraystretch}{0.8}
\setlength{\tabcolsep}{4pt}
\begin{tabular}{|p{3cm}|p{9cm}|}
\hline
\textbf{Élément} & \textbf{Description} \\
\hline
Cas d'utilisation & Traiter un ticket et mettre à jour le statut \\
\hline
Acteurs & Technicien (secondaires : Client, Administrateur notifiés) \\
\hline
Précondition & Le technicien est authentifié. Le ticket lui est \textbf{assigné} (\texttt{assignedTo}). \\
\hline
Postcondition & Le statut est mis à jour (\texttt{open}, \texttt{in\_progress} ou \texttt{resolved}) ; historique \texttt{status\_changed}. Client et admins notifiés. \\
\hline
Scénario nominal &
1. Le technicien accède à \texttt{/technician/tickets}.\\
& 2. \texttt{GET /tickets} retourne les tickets qui lui sont assignés.\\
& 3. Il change le statut via la liste déroulante (À traiter / En cours / Résolu).\\
& 4. \texttt{PATCH /tickets/:id} avec \texttt{\{ status \}} uniquement.\\
& 5. \texttt{handleTechnicianUpdate()} vérifie l'assignation et enregistre le changement.\\
& 6. Notifications temps réel et persistantes émises.\\
\hline
Scénarios alternatifs &
• Ticket non assigné à ce technicien (403).\\
& • Tentative de clôture \texttt{closed} par le technicien (403 — réservé à l'admin).\\
& • Tentative de modifier l'affectation (403).\\
& • Session expirée (401).\\
\hline
\end{tabular}
\caption{Description textuelle du cas d'utilisation : Traiter un ticket}
\label{tab:cu-process-ticket}
\end{table}

La figure~\ref{fig:seq-process-ticket} illustre ce scénario.

\begin{figure}[htbp]
    \centering
    \includegraphics[width=0.85\textwidth]{img/seq-traitement-ticket-statut.png}
    \caption{Diagramme de séquence -- Traitement et mise à jour du statut}
    \label{fig:seq-process-ticket}
\end{figure}

\subsubsection{Développement et intégration de l'interface de traitement de ticket}
L'espace technicien affiche les tickets assignés avec filtres. Le changement de statut se fait directement depuis la liste (sans module de commentaires). Le client reçoit une notification à chaque évolution du statut.


% ============================================================
% US06 — Gestion des utilisateurs
% ============================================================
\subsubsection{Description textuelle du cas d'utilisation « Gérer les utilisateurs »}

\begin{table}[H]
\centering
\footnotesize
\renewcommand{\arraystretch}{0.8}
\setlength{\tabcolsep}{4pt}
\begin{tabular}{|p{3cm}|p{9cm}|}
\hline
\textbf{Élément} & \textbf{Description} \\
\hline
Cas d'utilisation & Gérer les utilisateurs (CRUD) \\
\hline
Acteurs & Administrateur \\
\hline
Précondition & L'administrateur est authentifié (rôle \textit{admin}). Routes protégées par \texttt{RolesGuard}. \\
\hline
Postcondition & Les opérations CRUD sur la collection \texttt{users} sont effectuées ; la liste admin est rafraîchie. \\
\hline
Scénario nominal &
1. L'administrateur accède à \texttt{/admin} (section utilisateurs).\\
& 2. \texttt{GET /users} avec filtres \texttt{role}, \texttt{search}, pagination.\\
& 3. Il crée (\texttt{POST /users}), modifie (\texttt{PATCH /users/:id}) ou supprime (\texttt{DELETE /users/:id}).\\
& 4. Mot de passe hashé bcrypt à la création/modification.\\
\hline
Scénarios alternatifs &
• Email déjà utilisé (409 — \texttt{ConflictException}).\\
& • Suppression d'un technicien avec tickets assignés (409).\\
& • Session expirée (401).\\
\hline
\end{tabular}
\caption{Description textuelle du cas d'utilisation : Gérer les utilisateurs}
\label{tab:cu-manage-users}
\end{table}

La figure~\ref{fig:seq-manage-users} illustre ce scénario.

\begin{figure}[htbp]
    \centering
    \includegraphics[width=0.85\textwidth]{img/seq-gestion-utilisateurs.png}
    \caption{Diagramme de séquence -- Gestion des utilisateurs}
    \label{fig:seq-manage-users}
\end{figure}

\subsubsection{Développement et intégration de l'interface de gestion des utilisateurs}
L'interface admin centralise la liste paginée des utilisateurs avec recherche et filtre par rôle. Les formulaires \texttt{AddUserForm} et \texttt{EditUserModal} permettent le CRUD complet (client, technicien, admin).


% ============================================================
% US07 — Tableau de bord et statistiques
% ============================================================
\subsubsection{Description textuelle du cas d'utilisation « Consulter le tableau de bord »}

\begin{table}[H]
\centering
\footnotesize
\renewcommand{\arraystretch}{0.8}
\setlength{\tabcolsep}{4pt}
\begin{tabular}{|p{3cm}|p{9cm}|}
\hline
\textbf{Élément} & \textbf{Description} \\
\hline
Cas d'utilisation & Consulter le tableau de bord et les statistiques \\
\hline
Acteurs & Administrateur \\
\hline
Précondition & L'administrateur est authentifié (rôle \textit{admin}). \\
\hline
Postcondition & Les indicateurs agrégés sont affichés : utilisateurs par rôle, tickets par statut, priorité et technicien. \\
\hline
Scénario nominal &
1. L'administrateur accède au dashboard \texttt{/admin}.\\
& 2. \texttt{GET /users/stats} retourne \texttt{\{ total, byRole \}}.\\
& 3. \texttt{GET /tickets/stats} retourne \texttt{\{ byStatus, byPriority, byTechnician, total \}}.\\
& 4. Affichage via \texttt{StatsCards} et graphiques Chart.js (Doughnut, Bar).\\
\hline
Scénarios alternatifs &
• Base vide (compteurs à zéro).\\
& • Erreur chargement API (401/500).\\
\hline
\end{tabular}
\caption{Description textuelle du cas d'utilisation : Consulter le tableau de bord}
\label{tab:cu-dashboard}
\end{table}

La figure~\ref{fig:seq-dashboard} illustre ce scénario.

\begin{figure}[htbp]
    \centering
    \includegraphics[width=0.85\textwidth]{img/seq-tableau-de-bord.png}
    \caption{Diagramme de séquence -- Tableau de bord et statistiques}
    \label{fig:seq-dashboard}
\end{figure}

\subsubsection{Développement et intégration de l'interface du tableau de bord}
Le tableau de bord admin agrège les statistiques utilisateurs et tickets via deux endpoints dédiés. Les graphiques Chart.js visualisent la répartition par statut, par priorité et le volume par technicien assigné.


% ============================================================
% US08 — Notifications persistantes et alertes temps réel
% ============================================================
\subsubsection{Description textuelle du cas d'utilisation « Recevoir des notifications »}

\begin{table}[H]
\centering
\footnotesize
\renewcommand{\arraystretch}{0.8}
\setlength{\tabcolsep}{4pt}
\begin{tabular}{|p{3cm}|p{9cm}|}
\hline
\textbf{Élément} & \textbf{Description} \\
\hline
Cas d'utilisation & Recevoir des notifications et alertes temps réel \\
\hline
Acteurs & Client, Technicien, Administrateur \\
\hline
Précondition & L'acteur est authentifié. Connexion Socket.IO établie sur le namespace \texttt{/notifications} (JWT). \\
\hline
Postcondition & L'acteur reçoit une alerte temps réel (toast) et peut consulter l'historique persistant via \texttt{NotificationsBell}. \\
\hline
Scénario nominal &
1. À la connexion, le frontend ouvre Socket.IO avec le JWT.\\
& 2. Un événement ticket se produit (\texttt{ticket.created}, \texttt{ticket.assigned}, \texttt{ticket.status\_changed}).\\
& 3. \texttt{TicketsNotificationsGateway} émet l'événement WS vers les rooms ciblées.\\
& 4. \texttt{NotificationsService} persiste la notification en MongoDB.\\
& 5. \texttt{NotificationsBell} affiche le compteur et la liste (\texttt{GET /notifications}).\\
\hline
Scénarios alternatifs &
• Perte WebSocket (reconnexion automatique Socket.IO).\\
& • Notifications non lues consultables à la reconnexion (persistance).\\
& • Token expiré (déconnexion et redirection \texttt{/login}).\\
\hline
\end{tabular}
\caption{Description textuelle du cas d'utilisation : Recevoir des notifications}
\label{tab:cu-notifications}
\end{table}

La figure~\ref{fig:seq-notifications} illustre ce scénario.

\begin{figure}[htbp]
    \centering
    \includegraphics[width=0.85\textwidth]{img/seq-notifications.png}
    \caption{Diagramme de séquence -- Notifications temps réel (WebSocket)}
    \label{fig:seq-notifications}
\end{figure}

\subsubsection{Développement et intégration du système de notifications}
Le système combine Socket.IO (alertes instantanées + rafraîchissement des listes) et persistance MongoDB (cloche \texttt{NotificationsBell}, marquage lu/non lu via \texttt{PATCH /notifications/read-all}).


% ============================================================
% US09 — Historique et traçabilité
% ============================================================
\subsubsection{Description textuelle du cas d'utilisation « Consulter l'historique »}

\begin{table}[H]
\centering
\footnotesize
\renewcommand{\arraystretch}{0.8}
\setlength{\tabcolsep}{4pt}
\begin{tabular}{|p{3cm}|p{9cm}|}
\hline
\textbf{Élément} & \textbf{Description} \\
\hline
Cas d'utilisation & Consulter l'historique et la traçabilité d'un ticket \\
\hline
Acteurs & Client, Technicien, Administrateur \\
\hline
Précondition & L'acteur est authentifié et a accès au ticket (propriétaire, assigné ou admin). \\
\hline
Postcondition & La chronologie des actions (\texttt{history[]}) est affichée dans la modal de détail. \\
\hline
Scénario nominal &
1. L'acteur sélectionne un ticket dans sa liste.\\
& 2. \texttt{GET /tickets/:id} retourne le ticket avec le tableau \texttt{history[]}.\\
& 3. \texttt{TicketDetailModal} affiche la timeline horodatée.\\
& 4. Chaque entrée indique l'acteur, l'action (\texttt{created}, \texttt{updated}, \texttt{status\_changed}, \texttt{assigned\_changed}) et les valeurs avant/après.\\
\hline
Scénarios alternatifs &
• Historique vide (ticket sans entrée — cas rare).\\
& • Ticket introuvable ou accès refusé (403/404).\\
& • Session expirée (401).\\
\hline
\end{tabular}
\caption{Description textuelle du cas d'utilisation : Consulter l'historique}
\label{tab:cu-history}
\end{table}

La figure~\ref{fig:seq-history} illustre ce scénario.

\begin{figure}[htbp]
    \centering
    \includegraphics[width=0.85\textwidth]{img/seq-historique.png}
    \caption{Diagramme de séquence -- Historique et traçabilité}
    \label{fig:seq-history}
\end{figure}

\subsubsection{Développement et intégration de l'interface d'historique}
L'historique est embarqué dans le document ticket MongoDB (\texttt{TicketHistoryEntry[]}). Il est rendu dans \texttt{TicketDetailModal} sous forme de liste chronologique, sans onglet séparé ni filtre avancé dans la version Sprint~1.
